\documentclass[11pt]{article}
\usepackage[margin=1in]{geometry}
\usepackage{amsmath,amssymb}
\usepackage{bm}
\usepackage{booktabs}
\usepackage{siunitx}
\usepackage{xcolor}
\usepackage[colorlinks=true,linkcolor=blue,urlcolor=blue]{hyperref}
\usepackage{caption}
\captionsetup{font=small}

\newcommand{\R}{\mathbb{R}}
\renewcommand{\d}{\mathrm{d}}

\title{\vspace{-2em}Aitken acceleration in Norma: two formulations}
\author{Alejandro Mota \\ \small Sandia National Laboratories}
\date{\today}

\begin{document}
\maketitle

\begin{abstract}
This note records the two Aitken-acceleration formulations implemented in
Norma's non-overlapping Schwarz alternating solver, together with the
displacement-form kinematic recovery they use. The purpose is to fix the exact
formulas and index conventions so they can be reproduced in Sierra/SM. Both
formulations are exposed through the multi-domain controller: the recursive
Irons--Tuck variant (\texttt{relaxation: aitken}) and the non-recursive secant
variant (\texttt{relaxation: aitken secant}). They are mathematically
equivalent in exact arithmetic but differ in their initialization bookkeeping
and in how they treat the interface velocity and acceleration.
\end{abstract}

\section{Setting and fixed-point reformulation}

Consider a domain decomposition of $\Omega$ into two non-overlapping subdomains
$\Omega_1$ and $\Omega_2$ sharing the interface $\Gamma$. Norma drives the
coupling \emph{multiplicatively} (Gauss--Seidel ordering of the subdomain
solves) with a strongly imposed Dirichlet condition on the Dirichlet side of
the interface and a traction (Neumann) condition on the other. Let $\bm{g}$
denote the interface unknown, i.e.\ the displacement datum imposed on the
interface. Following Sambataro and Tezaur~\cite{SambataroTezaur} we write the
multiplicative Schwarz sweep as a fixed-point map
\begin{equation}
    \bm{g} \;\mapsto\; T(\bm{g}),
    \label{eq:T}
\end{equation}
where $T(\bm{g})$ is obtained by solving the subdomains in order with the
current interface data and projecting the resulting trace back onto $\Gamma$;
the converged interface state satisfies $\bm{g}^\star = T(\bm{g}^\star)$.

Define the \emph{interface jump} (the fixed-point residual)
\begin{equation}
    \bm{r}^{(n)} \;:=\; T\!\left(\bm{g}^{(n)}\right) - \bm{g}^{(n)} .
    \label{eq:residual}
\end{equation}
In the Norma implementation (\texttt{src/schwarz.jl}) the two terms map to the
code variables as
\begin{equation*}
    \bm{g}^{(n)} \;=\; \texttt{lambda\_disp} \ \text{(from iteration $n{-}1$)},
    \qquad
    T\!\left(\bm{g}^{(n)}\right) \;=\; \texttt{interp\_disp} ,
\end{equation*}
so that $\bm{r}^{(n)} = \texttt{interp\_disp} - \texttt{lambda\_disp}$. This is
the displacement mismatch the relaxation acts on; it is non-degenerate because
\texttt{interp\_disp} is the freshly projected neighbor trace, whereas the
strongly imposed continuity makes the \emph{post-imposition} interface residual
identically zero on the Dirichlet side.

The relaxed update common to all variants is
\begin{equation}
    \bm{g}^{(n+1)} \;=\; \bm{g}^{(n)} + \omega^{(n)} \, \bm{r}^{(n)}
    \;=\; \omega^{(n)} \, T\!\left(\bm{g}^{(n)}\right)
        + \bigl(1 - \omega^{(n)}\bigr)\,\bm{g}^{(n)} ,
    \label{eq:update}
\end{equation}
with relaxation factor $\omega^{(n)}$. Classical relaxation uses a fixed
$\omega^{(n)} \equiv \omega$; the Aitken variants choose $\omega^{(n)}$
dynamically.

\section{Version 1: recursive Irons--Tuck Aitken (\texttt{relaxation: aitken})}

With $\bm{\delta}^{(n)} := \bm{r}^{(n)} - \bm{r}^{(n-1)}$, the Irons--Tuck
recursion~\cite{IronsTuck} sets
\begin{equation}
    \omega^{(n)} \;=\;
    -\,\omega^{(n-1)}\,
    \frac{\bm{r}^{(n-1)} \cdot \bm{\delta}^{(n)}}
         {\bigl\lVert \bm{\delta}^{(n)} \bigr\rVert^{2}} .
    \label{eq:ironstuck}
\end{equation}
The factor carries the damping history through the multiplicative
$\omega^{(n-1)}$, which requires the iterate to have advanced by exactly
$\omega^{(n-1)}\bm{r}^{(n-1)}$ for the recursion to remain consistent (see
Section~\ref{sec:equiv}).

\section{Version 2: non-recursive secant Aitken (\texttt{relaxation: aitken secant})}

This is the form derived in the original references,
Sambataro--Tezaur~\cite{SambataroTezaur} (Eq.~9) and
Deparis--Discacciati--Quarteroni~\cite{Deparis}. With
$\bm{d}^{(n)} := \bm{g}^{(n)} - \bm{g}^{(n-1)}$ (the change in the interface
iterate) and $\bm{\delta}^{(n)}$ as above,
\begin{equation}
    \omega^{(n)} \;=\;
    -\,\frac{\bm{d}^{(n)} \cdot \bm{\delta}^{(n)}}
            {\bigl\lVert \bm{\delta}^{(n)} \bigr\rVert^{2}}
    \;=\; \arg\min_{\omega}\,
          \bigl\lVert \bm{d}^{(n)} + \omega\,\bm{\delta}^{(n)} \bigr\rVert^{2} .
    \label{eq:secant}
\end{equation}
The numerator uses the change of the iterate $\bm{d}^{(n)}$, not the previous
residual, and there is no recursive $\omega^{(n-1)}$ factor. Norma forms
$\bm{d}^{(n)} = \bm{g}^{(n)} - \bm{g}^{(n-1)}$ directly from the two stored
interface iterates rather than from the recursion.

\section{Relationship between the two forms}
\label{sec:equiv}

The update~\eqref{eq:update} gives
$\bm{g}^{(n)} = \bm{g}^{(n-1)} + \omega^{(n-1)}\bm{r}^{(n-1)}$, hence
\begin{equation}
    \bm{d}^{(n)} \;=\; \bm{g}^{(n)} - \bm{g}^{(n-1)}
                 \;=\; \omega^{(n-1)}\,\bm{r}^{(n-1)} .
    \label{eq:dlink}
\end{equation}
Substituting~\eqref{eq:dlink} into the secant factor~\eqref{eq:secant}
reproduces the Irons--Tuck factor~\eqref{eq:ironstuck} exactly. The two forms
therefore coincide \emph{whenever the interface iterate advances by exactly}
$\omega^{(n-1)}\bm{r}^{(n-1)}$. They diverge only where that identity is broken
by the initialization and warm-up bookkeeping (Section~\ref{sec:init}): the
secant form reads the actual stored iterates $\bm{g}^{(n)},\bm{g}^{(n-1)}$ and
is robust to it, while the recursion silently inherits any inconsistency in the
stored $\omega^{(n-1)}$ and $\bm{r}^{(n-1)}$.

\section{An equivalent map-output form of the secant factor}
\label{sec:mapoutput}

An alternative statement of the acceleration, proposed independently by
D.~Koliesnikova, works with the differences of the \emph{map outputs}
$T(\bm{g}^{(n)})$ rather than of the iterates. Define
\begin{equation}
    \beta^{(n)} \;=\;
    \frac{\bigl(T(\bm{g}^{(n)}) - T(\bm{g}^{(n-1)})\bigr)\cdot\bm{\delta}^{(n)}}
         {\bigl\lVert \bm{\delta}^{(n)} \bigr\rVert^{2}},
    \qquad
    \omega^{(n)} \;=\; 1 - \beta^{(n)},
    \label{eq:mapoutput}
\end{equation}
followed by the same relaxed update~\eqref{eq:update}, equivalently written
$\bm{g}^{(n+1)} = T(\bm{g}^{(n)}) - \beta^{(n)}\bm{r}^{(n)}$. Since
$T(\bm{g}) = \bm{g} + \bm{r}$, the map-output difference decomposes as
\begin{equation}
    T(\bm{g}^{(n)}) - T(\bm{g}^{(n-1)})
    \;=\; \bm{d}^{(n)} + \bm{\delta}^{(n)},
    \label{eq:Tdecomp}
\end{equation}
so that
\begin{equation*}
    \omega^{(n)} \;=\; 1 -
    \frac{\bigl(\bm{d}^{(n)} + \bm{\delta}^{(n)}\bigr)\cdot\bm{\delta}^{(n)}}
         {\lVert \bm{\delta}^{(n)} \rVert^{2}}
    \;=\; -\,\frac{\bm{d}^{(n)}\cdot\bm{\delta}^{(n)}}
                  {\lVert \bm{\delta}^{(n)} \rVert^{2}} ,
\end{equation*}
which is exactly the secant factor~\eqref{eq:secant}. The map-output
form~\eqref{eq:mapoutput} is therefore algebraically identical to
Version~2.

\paragraph{Notation.} The original statement of~\eqref{eq:mapoutput}
writes the fixed-point residual itself as $\delta g^{n}$, so the
quantity $\delta g^{n} - \delta g^{n-1}$ appearing there in the
numerator and denominator is this note's residual difference
$\bm{\delta}^{(n)}$ --- the shared letter $\delta$ denotes different
quantities in the two notations. It also indexes the sweep
\emph{output} ($g^{n} = T(g^{n-1})$), whereas this note evaluates
$T$ and $\bm{r}$ at a single iterate level; the residual definitions
coincide after shifting her iterate indices by one.
Table~\ref{tab:translation} gives the correspondence.

\begin{table}[ht!]
\centering
\small
\begin{tabular}{@{}lll@{}}
\toprule
Original statement & This note & Meaning \\
\midrule
$\delta g^{n} = T(g^{n}) - g^{n-1}$ &
$\bm{r}^{(n)} = T(\bm{g}^{(n)}) - \bm{g}^{(n)}$ &
fixed-point residual \\
$\delta g^{n} - \delta g^{n-1}$ &
$\bm{\delta}^{(n)} = \bm{r}^{(n)} - \bm{r}^{(n-1)}$ &
residual difference \\
$T(g^{n}) - T(g^{n-1})$ &
$\bm{d}^{(n)} + \bm{\delta}^{(n)}$ &
map-output difference \\
\bottomrule
\end{tabular}
\caption{Correspondence between the notation of the original statement
of~\eqref{eq:mapoutput} and that of this note. The mixed iterate
indices in the first row reflect the output-based indexing
$g^{n} = T(g^{n-1})$; both rows denote the same residual once the
iterate indices are aligned.}
\label{tab:translation}
\end{table}

The form~\eqref{eq:mapoutput} is in particular \emph{not} the Irons--Tuck
recursion~\eqref{eq:ironstuck}: like the secant form, it is non-recursive,
reads only stored data, and inherits the robustness discussed in
Section~\ref{sec:equiv}. That the same factor is reached from two
independent derivations is a useful cross-check on both.

Two implementation remarks. First, the storage requirements are the same
but the stored quantities differ: \eqref{eq:mapoutput} retains the two map
outputs (or one output and the residual history), whereas the
implementation retains the iterate $\bm{g}^{(n-1)}$ and the residual
$\bm{r}^{(n-1)}$. Second, computing $\omega^{(n)} = 1 - \beta^{(n)}$
incurs a cancellation when $\beta^{(n)} \to 1$ (strong damping,
$\omega^{(n)}$ small) that amplifies the \emph{relative} error of
$\omega^{(n)}$, although the absolute error --- which is what enters the
update --- is unaffected; the direct form~\eqref{eq:secant} avoids even
this. The considerations that distinguish the variants in practice are
unchanged by the choice of form: the warm-up and safeguard bookkeeping
(Sections~\ref{sec:init}--\ref{sec:fields}) and, most consequentially for
dynamics, the choice of which interface variable to relax
(Section~\ref{sec:fields}), on which~\eqref{eq:mapoutput} is silent.

\section{Initialization}
\label{sec:init}

Let $\omega^{(1)}$ be a user-supplied initial factor (Norma:
\texttt{relaxation parameter}, default $1$) and $N_0$ a warm-up count (Norma:
\texttt{aitken N0 parameter}, default $1$). For $1 \le n < N_0$ the factor is
held at $\omega^{(n)} = \omega^{(1)}$; the first dynamic factor is computed at
$n = N_0$, once two residuals (Version~1) and additionally two iterates
(Version~2) are on record.

\section{Numerical safeguards}

A single safeguard is retained: a floor on
$\lVert\bm{\delta}^{(n)}\rVert^{2}$ ($10^{-20}$) that guards the division
in~\eqref{eq:ironstuck} and~\eqref{eq:secant} when successive residuals are
essentially identical; in that case the previous factor is reused (Version~1)
or $\omega^{(1)}$ is used (Version~2). No value clamp is applied to
$\omega^{(n)}$: clamping caps the very large or near-zero steps that produce
the acceleration and was found to slow convergence, so it is omitted.

\section{Which variable to relax: kinematic consistency}
\label{sec:fields}

For dynamics the interface carries displacement $\bm{u}$, velocity $\bm{v}$,
and acceleration $\bm{a}$. Under Newmark integration with a \emph{fixed}
step history $(\bm{u}_n,\bm{v}_n,\bm{a}_n)$ and parameters $\beta,\gamma,\Delta
t$, the end-of-step kinematics lie on a one-parameter family,
\begin{equation}
    \bm{u}_{n+1} = \bm{u}_{\mathrm{pre}} + \beta\Delta t^2\,\bm{a}_{n+1},
    \qquad
    \bm{v}_{n+1} = \bm{v}_{\mathrm{pre}} + \gamma\Delta t\,\bm{a}_{n+1},
    \label{eq:newmark}
\end{equation}
with predictors
$\bm{u}_{\mathrm{pre}} = \bm{u}_n + \Delta t\,\bm{v}_n
 + (\tfrac12 - \beta)\Delta t^2\,\bm{a}_n$ and
$\bm{v}_{\mathrm{pre}} = \bm{v}_n + (1-\gamma)\Delta t\,\bm{a}_n$.
Thus only one of $(\bm{u},\bm{v},\bm{a})$ is independent: the integrator's
primary solve variable. The relaxation should act on that variable, with the
remaining two recovered from~\eqref{eq:newmark}.

\paragraph{Norma (displacement-form Newmark).}
The solver unknown is $\bm{u}$ (\texttt{src/time\_integrator.jl}, \texttt{correct}).
The secant variant therefore relaxes \emph{displacement only},
$\bm{u}^{(n+1)} = \bm{g}^{(n+1)}$ from~\eqref{eq:update}, and recovers
\begin{equation}
    \bm{a}^{(n+1)} = \frac{\bm{u}^{(n+1)} - \bm{u}_{\mathrm{pre}}}{\beta\Delta t^2},
    \qquad
    \bm{v}^{(n+1)} = \bm{v}_{\mathrm{pre}} + \gamma\Delta t\,\bm{a}^{(n+1)} .
    \label{eq:recover}
\end{equation}
The legacy \texttt{relaxation: aitken} variant instead relaxes
$\bm{u}$, $\bm{v}$, and $\bm{a}$ independently with the same $\omega^{(n)}$. A
$\omega$-blend of two kinematically consistent triples is again consistent, so
the two treatments agree when the incoming triple lies on~\eqref{eq:newmark};
they differ slightly otherwise because Norma interpolates the three interface
histories independently in time. Recovering from a single relaxed variable
keeps the imposed interface state exactly on the integrator manifold.

\paragraph{Sierra/SM (acceleration-form).}
In an acceleration-form integrator the primary unknown is $\bm{a}$. The same
principle then prescribes relaxing the \emph{acceleration} and recovering
$\bm{u}$ and $\bm{v}$ by inverting~\eqref{eq:newmark}, i.e.\ replacing the role
of $\bm{u}$ in~\eqref{eq:recover} by $\bm{a}$. This is the relevant choice for
reproduction in Sierra/SM.

\section{Summary for reproduction}

\begin{table}[ht!]
\centering
\small
\begin{tabular}{@{}llll@{}}
\toprule
Formulation & Relaxation factor & Norma key & Variable relaxed \\
\midrule
Irons--Tuck (recursive) &
$\omega^{(n)} = -\omega^{(n-1)}\dfrac{\bm{r}^{(n-1)}\cdot\bm{\delta}^{(n)}}{\lVert\bm{\delta}^{(n)}\rVert^2}$ &
\texttt{aitken} & $\bm{u},\bm{v},\bm{a}$ (same $\omega$) \\[1.4em]
Secant (non-recursive) &
$\omega^{(n)} = -\dfrac{\bm{d}^{(n)}\cdot\bm{\delta}^{(n)}}{\lVert\bm{\delta}^{(n)}\rVert^2}$ &
\texttt{aitken secant} & primary unknown; recover rest \\
\bottomrule
\end{tabular}
\caption{The two Aitken formulations implemented in Norma. Here
$\bm{r}^{(n)} = T(\bm{g}^{(n)}) - \bm{g}^{(n)}$,
$\bm{\delta}^{(n)} = \bm{r}^{(n)} - \bm{r}^{(n-1)}$, and
$\bm{d}^{(n)} = \bm{g}^{(n)} - \bm{g}^{(n-1)}$. For Norma's
displacement-form Newmark the primary unknown is $\bm{u}$; for an
acceleration-form integrator (Sierra/SM) it is $\bm{a}$.}
\end{table}

On a representative dynamic non-overlapping Dirichlet--Neumann cantilever
problem the secant variant converged in \num{2} Schwarz iterations per time
step, versus \numrange{4}{5} for the recursive variant. The two are not
interchangeable there, however: because the displacement-only secant imposes a
recovered (rather than independently relaxed) interface velocity and
acceleration, its transient differs from the all-fields recursive variant. A
full-field comparison against the monolithic single-domain solution shows the
recursive variant within \SI{0.7}{\percent} while the secant variant is
\SI{8.8}{\percent} off; the agreement is close only at isolated points such as
the cantilever tip. The secant variant thus accelerates the displacement
fixed point at the cost of transient accuracy for Dirichlet--Neumann dynamics.
For the force-relaxed couplings (Robin--Robin, impedance) no kinematic recovery
is involved and all methods reproduce the monolithic solution equally well
(\SI{0.5}{\percent}--\SI{0.7}{\percent} on the same problem).

\begin{thebibliography}{9}
\bibitem{SambataroTezaur} G.~Sambataro and I.~Tezaur,
\emph{On the role of relaxation and acceleration in the non-overlapping Schwarz
alternating method for coupling}, arXiv:2603.17143.
\bibitem{Deparis} S.~Deparis, M.~Discacciati, and A.~Quarteroni,
\emph{A domain decomposition framework for fluid-structure interaction
problems}, MOX Report 45, Politecnico di Milano.
\bibitem{IronsTuck} B.~M.~Irons and R.~C.~Tuck,
\emph{A version of the Aitken accelerator for computer iteration},
Int.\ J.\ Numer.\ Methods Eng.\ \textbf{1} (1969) 275--277.
\end{thebibliography}

\end{document}
